% ============================================================================
% LANGENTWURF: Repaso - Mi ropa preferida
% Klasse 7 | Spanisch | Doppelstunde (90 Min.)
% ============================================================================
% Lehrwerk: Qué pasa 1, Unidad 5 (S. 88-101)
% Fokus: Wiederholung Kleidung, Farben, llevar, Einkaufsdialoge
% ============================================================================

\documentclass[11pt,a4paper]{article}

% ============================================================================
% PACKAGES
% ============================================================================
\usepackage[utf8]{inputenc}
\usepackage[T1]{fontenc}
\usepackage[ngerman]{babel}
\usepackage{geometry}
\usepackage{fancyhdr}
\usepackage{lastpage}
\usepackage{xcolor}
\usepackage{tcolorbox}
\usepackage{enumitem}
\usepackage{tabularx}
\usepackage{longtable}
\usepackage{array}
\usepackage{booktabs}
\usepackage{hyperref}
\usepackage{ifthen}
\usepackage{amssymb}

% ============================================================================
% KONFIGURATION
% ============================================================================

\newcommand{\plantier}{1}
\newcommand{\fachid}{spanisch}
\newcommand{\fach}{Spanisch}
\newcommand{\fachfarbe}{c41e3a}

% ============================================================================
% VARIABLEN
% ============================================================================

\newcommand{\jahrgangsstufe}{7}
\newcommand{\schulart}{Gesamtschule}
\newcommand{\stundenthema}{Repaso -- Mi ropa preferida}
\newcommand{\reihenthema}{Unidad 5: Mi ropa preferida}
\newcommand{\stundennummer}{Repaso}
\newcommand{\reihenstunden}{4}
\newcommand{\unterrichtsdatum}{\_\_\_\_\_\_\_\_\_\_}
\newcommand{\unterrichtszeit}{\_\_\_\_\_ -- \_\_\_\_\_}
\newcommand{\raum}{\_\_\_\_\_}

\newcommand{\lehrkraft}{N.N.}
\newcommand{\ausbildungsstand}{Lehrkraft}

\newcommand{\lerngruppengroesse}{24}
\newcommand{\maennlich}{12}
\newcommand{\weiblich}{12}
\newcommand{\divers}{0}

\newcommand{\gerniveau}{A1}
\newcommand{\lehrwerk}{Qué pasa 1}
\newcommand{\lehrwerkseite}{S. 88-101}

% ============================================================================
% LAYOUT
% ============================================================================
\geometry{left=2.5cm, right=2.5cm, top=2.5cm, bottom=2.5cm}
\definecolor{fach}{HTML}{\fachfarbe}
\definecolor{gray}{HTML}{666666}
\definecolor{lightgray}{HTML}{f5f5f5}
\definecolor{successgreen}{HTML}{2a7a4b}
\definecolor{warningorange}{HTML}{b35c00}

\tcbuselibrary{skins,breakable}

\newtcolorbox{infobox}[1][]{
    colback=lightgray,
    colframe=fach,
    fonttitle=\bfseries,
    title=#1,
    breakable
}

\newtcolorbox{scorebox}{
    colback=white,
    colframe=fach,
    fonttitle=\bfseries,
    title=Didaktik-Score,
    breakable
}

\pagestyle{fancy}
\fancyhf{}
\fancyhead[L]{\small\textcolor{gray}{\fach{} -- Jg. \jahrgangsstufe{} -- \stundenthema}}
\fancyhead[R]{\small\textcolor{gray}{Langentwurf Tier 1}}
\fancyfoot[C]{\small\thepage{} / \pageref{LastPage}}
\renewcommand{\headrulewidth}{0.4pt}
\renewcommand{\footrulewidth}{0pt}

% ============================================================================
% DOKUMENT
% ============================================================================
\begin{document}

% ============================================================================
% DECKBLATT
% ============================================================================
\begin{titlepage}
    \centering
    \vspace*{2cm}

    {\large\textcolor{gray}{\schulart}}\\[2cm]

    {\Huge\textcolor{fach}{\textbf{Langentwurf}}}\\[0.5cm]
    {\large Tier 1 -- Vollständig}\\[2cm]

    {\LARGE\textbf{\stundenthema}}\\[0.5cm]
    {\large im Rahmen der Unterrichtsreihe}\\[0.3cm]
    {\Large\textit{\reihenthema}}\\[2cm]

    \begin{tabular}{rl}
        \textbf{Fach:} & \fach{} (\gerniveau) \\[0.3cm]
        \textbf{Klasse:} & \jahrgangsstufe \\[0.3cm]
        \textbf{Lehrwerk:} & \lehrwerk, \lehrwerkseite \\[0.3cm]
        \textbf{Datum:} & \underline{\hspace{3cm}} \\[0.3cm]
        \textbf{Zeit:} & \underline{\hspace{3cm}} (Doppelstunde) \\[0.3cm]
        \textbf{Raum:} & \underline{\hspace{3cm}} \\[0.3cm]
    \end{tabular}

    \vfill

    {\large Lehrkraft: \lehrkraft}\\[0.3cm]
    {\normalsize\textcolor{gray}{\ausbildungsstand}}\\[1cm]

\end{titlepage}

% ============================================================================
% INHALTSVERZEICHNIS
% ============================================================================
\tableofcontents
\newpage

% ============================================================================
% 1. BEDINGUNGSANALYSE
% ============================================================================
\section{Bedingungsanalyse}

\subsection{Lerngruppe}
\begin{tabular}{ll}
    Kursgröße: & \lerngruppengroesse{} SuS (\maennlich{} m / \weiblich{} w / \divers{} d) \\
    Jahrgangsstufe: & \jahrgangsstufe \\
    Sprachniveau: & \gerniveau{} (GER) \\
    Fremdsprache: & 2. Fremdsprache (1. Lernjahr) \\
\end{tabular}

\subsection{Vorkenntnisse}
Die SuS haben aus Unidad 5 bereits gelernt:
\begin{itemize}
    \item Kleidungsvokabeln (la camiseta, los pantalones, la falda, el vestido, etc.)
    \item Farben und Farbadjektive (rojo/a, azul, verde, etc.)
    \item Verb \textit{llevar} (Yo llevo una camiseta roja.)
    \item Verb \textit{ir} (Vamos de compras.)
    \item Modalverben \textit{querer}, \textit{poder}, \textit{tener que}
    \item Spaltungsverb \textit{costar} (o$\rightarrow$ue)
    \item Zahlen 100-1000 (für Preise)
\end{itemize}

\textbf{Fokus dieser Repaso-Stunde:} Festigung und kommunikative Anwendung des Gelernten.

\subsection{Differenzierung (intern)}

\textbf{WICHTIG:} Keine Sternchen auf Schülermaterial!

\begin{tabular}{|l|p{4cm}|p{4cm}|p{4cm}|}
\hline
& \textbf{[B] Basis} & \textbf{[S] Standard} & \textbf{[E] Erweiterung} \\
\hline
\textbf{Ziel} & Berufsreife & Mittlere Reife & Gymnasium \\
\hline
\textbf{Vokabeln} & 10 Kernvokabeln rezeptiv & 15 Vokabeln produktiv & 20+ Vokabeln + freie Anwendung \\
\hline
\textbf{Grammatik} & \textit{llevar} + Lückentext & Farbadjektive korrekt & Einkaufsdialog frei gestalten \\
\hline
\textbf{Produktion} & Sätze nach Muster & Geführter Dialog & Freie Modenschau-Präsentation \\
\hline
\end{tabular}

\subsection{Besonderheiten}
\begin{itemize}
    \item \textbf{LRS:} 2-3 SuS (Zeitzugabe, Vokabelliste mit Bildern)
    \item \textbf{DaZ:} 1-2 SuS (Farbvokabeln visuell unterstützen)
    \item \textbf{Leistungsstark:} 3-4 SuS (Modenschau moderieren, Expertenrolle)
\end{itemize}

% ============================================================================
% 2. SACHANALYSE
% ============================================================================
\section{Sachanalyse}

\subsection{Sprachliche Strukturen}

\subsubsection{Wortschatz: Kleidung (la ropa)}

\textbf{Kernvokabular:}
\begin{center}
\begin{tabular}{ll|ll}
\textbf{Spanisch} & \textbf{Deutsch} & \textbf{Spanisch} & \textbf{Deutsch} \\
\hline
la camiseta & das T-Shirt & el jersey & der Pullover \\
la camisa & das Hemd & la chaqueta & die Jacke \\
los pantalones & die Hose & la falda & der Rock \\
los vaqueros & die Jeans & el vestido & das Kleid \\
los zapatos & die Schuhe & las zapatillas & die Turnschuhe \\
las botas & die Stiefel & el cinturón & der Gürtel \\
\end{tabular}
\end{center}

\textbf{Farben (los colores):}
\begin{center}
\begin{tabular}{ll|ll|ll}
\textbf{Spanisch} & \textbf{Deutsch} & \textbf{Spanisch} & \textbf{Deutsch} & \textbf{Spanisch} & \textbf{Deutsch} \\
\hline
rojo/a & rot & azul & blau & verde & grün \\
amarillo/a & gelb & negro/a & schwarz & blanco/a & weiß \\
marrón & braun & gris & grau & rosa & rosa \\
naranja & orange & violeta & violett & & \\
\end{tabular}
\end{center}

\subsubsection{Grammatik: Wiederholung}

\textbf{1. Verb \textit{llevar} (tragen):}
\begin{center}
\begin{tabular}{ll}
yo llevo & nosotros/as llevamos \\
tú llevas & vosotros/as lleváis \\
él/ella/usted lleva & ellos/ellas/ustedes llevan \\
\end{tabular}
\end{center}

\textbf{2. Farbadjektive -- Kongruenz:}
\begin{itemize}
    \item 4-Endungen: rojo, roja, rojos, rojas
    \item 2-Endungen: azul, azules / verde, verdes / gris, grises
    \item Unveränderlich: rosa, naranja, violeta
\end{itemize}

\textbf{Beispiele:}
\begin{itemize}
    \item \textit{La camiseta \textbf{roja}} (f. Sg.)
    \item \textit{Los pantalones \textbf{azules}} (m. Pl.)
    \item \textit{La falda \textbf{rosa}} (unveränderlich)
\end{itemize}

\textbf{3. Verb \textit{costar} (o$\rightarrow$ue):}
\begin{itemize}
    \item Singular: \textit{¿Cuánto \textbf{cuesta} la camiseta?} -- \textit{Cuesta 15 euros.}
    \item Plural: \textit{¿Cuánto \textbf{cuestan} los zapatos?} -- \textit{Cuestan 45 euros.}
\end{itemize}

\subsection{Kontrastive Betrachtung (Deutsch $\leftrightarrow$ Spanisch)}
\begin{itemize}
    \item \textbf{Positiver Transfer:} Farbwörter nach Substantiv (wie im Französischen)
    \item \textbf{Interferenz:} Deutsche Schüler vergessen die Genus-Kongruenz (*la camiseta rojo)
    \item \textbf{Ausnahme:} Unveränderliche Farben (rosa, naranja) -- keine deutschen Äquivalente
\end{itemize}

\subsection{Kulturelle Aspekte}
\begin{itemize}
    \item \textbf{Schuluniformen:} In vielen lateinamerikanischen Ländern Pflicht, in Spanien selten
    \item \textbf{Mode in Spanien:} Zara, Mango, Desigual -- spanische Modemarken
    \item \textbf{Feria de Abril:} Traditionelle Kleidung (traje de flamenca) in Sevilla
\end{itemize}

% ============================================================================
% 3. DIDAKTISCHE ANALYSE
% ============================================================================
\section{Didaktische Analyse}

\subsection{Stellung in der Sequenz}
Dies ist die \textbf{Repaso-Doppelstunde} der Sequenz ``\reihenthema'' (nach 3-4 Einführungsstunden).

\textbf{Sequenzübersicht:}
\begin{enumerate}
    \item DS 1: Kleidungsvokabeln + \textit{llevar}
    \item DS 2: Farben + Adjektivkongruenz
    \item DS 3: Einkaufen + \textit{costar} + Zahlen 100-1000
    \item \textbf{DS 4 (diese Stunde): Repaso -- Wiederholung und Anwendung}
\end{enumerate}

\subsection{Kompetenzziele (Rahmenplan MV)}
\begin{itemize}
    \item \textbf{Kommunikativ (Sprechen):} Kleidung beschreiben, Einkaufsgespräche führen
    \item \textbf{Sprachlich (Wortschatz):} Kleidungs- und Farbvokabular sicher anwenden
    \item \textbf{Sprachlich (Grammatik):} Farbadjektive korrekt angleichen
    \item \textbf{Methodisch:} Vokabelwiederholung durch spielerische Übungen
\end{itemize}

\subsection{Legitimation}
\textbf{Rahmenplan Spanisch MV (Kl. 7-10), Kompetenzbereich Sprechen:}
\begin{quote}
``Die Schülerinnen und Schüler können in einfachen Sätzen Personen und Gegenstände beschreiben und in Alltagssituationen kommunizieren.''
\end{quote}

Das Thema Kleidung bietet einen authentischen, alltagsrelevanten Kontext für die Anwendung der Zielstrukturen.

% ============================================================================
% 4. STUNDENZIELE
% ============================================================================
\section{Stundenziele}

\subsection{Grobziel}
Die SuS festigen und erweitern ihre \textbf{kommunikative Kompetenz} im Bereich Sprechen, indem sie Kleidung beschreiben und einfache Einkaufsgespräche führen können.

\subsection{Feinziele (operationalisiert)}
\begin{tabular}{|c|p{8cm}|c|c|}
\hline
\textbf{Nr.} & \textbf{Die SuS können...} & \textbf{AFB} & \textbf{Diff.} \\
\hline
FZ1 & 12 Kleidungsvokabeln korrekt benennen & I & alle \\
\hline
FZ2 & Farbadjektive korrekt an das Substantiv angleichen & I & alle \\
\hline
FZ3 & mit \textit{llevar} Kleidung beschreiben & II & alle \\
\hline
FZ4 & Fragen mit \textit{¿De qué color...?} stellen und beantworten & II & [S]/[E] \\
\hline
FZ5 & einen Einkaufsdialog führen (\textit{costar}, Preise) & II & [S]/[E] \\
\hline
FZ6 & eine kurze Modenschau-Präsentation halten & III & [E] \\
\hline
\end{tabular}

% ============================================================================
% 5. METHODISCHE ANALYSE
% ============================================================================
\section{Methodische Analyse}

\subsection{Phasierung (Doppelstunde = 2 × 45 Min.)}

\textbf{1. Stunde: Vokabel-Wiederholung + Spiele}
\begin{itemize}
    \item Einstieg: ``Ich sehe was, was du nicht siehst'' (Farben)
    \item Wiederholung: Vokabel-Bingo (Kleidung)
    \item Übung I: Bildkarten-Zuordnung + Farbadjektive
    \item Übung II: Kettenübung ``Hoy llevo...''
\end{itemize}

\textbf{2. Stunde: Kommunikative Anwendung}
\begin{itemize}
    \item Aktivierung: Schnelles Vokabelquiz
    \item Anwendung: Einkaufsdialog (Partnerarbeit)
    \item Kreativ: Mini-Modenschau (Gruppenarbeit)
    \item Sicherung: Arbeitsblatt + Reflexion
\end{itemize}

\subsection{Medien}
\begin{itemize}
    \item Tafel/Whiteboard für Vokabel-Visualisierung
    \item Lehrwerk: Qué pasa 1, S. 88-101
    \item \textbf{AB-05-repaso-ropa.pdf} (Arbeitsblatt, max. 2 Seiten)
    \item \textbf{LB-05-repaso-ropa.pdf} (Lernblatt, max. 2 Seiten)
    \item Bingo-Karten (Kleidung)
    \item Bildkarten: Kleidungsstücke
    \item Optional: Echte Kleidungsstücke / Requisiten
\end{itemize}

\subsection{Sozialformen}
\begin{itemize}
    \item Plenum: Bingo, Kettenübung, Modenschau-Präsentation
    \item Partnerarbeit: Einkaufsdialoge, Beschreibungen
    \item Gruppenarbeit: Mini-Modenschau vorbereiten
    \item Einzelarbeit: Arbeitsblatt
\end{itemize}

% ============================================================================
% 6. VERLAUFSPLANUNG
% ============================================================================
\section{Verlaufsplanung}

\textbf{1. Stunde (45 Min.) -- Vokabel-Wiederholung}

\begin{longtable}{|p{1cm}|p{1.5cm}|p{4cm}|p{3cm}|p{2cm}|p{2cm}|}
\hline
\textbf{Zeit} & \textbf{Phase} & \textbf{Lehrerhandeln} & \textbf{Schülerhandeln} & \textbf{SF} & \textbf{Medien} \\
\hline
\endhead

0--5 & Einstieg & Begrüßung: \textit{¡Hola! Hoy hacemos un repaso.} Spiel: ``Veo, veo'' (Farben im Raum) & Raten Farben: \textit{¿Azul? ¿Rojo?} & Plenum & -- \\
\hline

5--15 & Wieder\-holung & Verteilt Bingo-Karten, ruft Kleidungsvokabeln auf Spanisch & Markieren auf Bingo-Karte, rufen \textit{¡Bingo!} & Plenum & Bingo-Karten \\
\hline

15--25 & Übung I & Zeigt Bildkarten, fragt: \textit{¿Qué es? ¿De qué color es?} & Antworten mit Farbadjektiv: \textit{Es una camiseta roja.} & Plenum & Bildkarten \\
\hline

25--35 & Übung II & Startet Kettenübung: \textit{Hoy llevo una camiseta azul.} & Beschreiben eigene Kleidung, reichen weiter & Plenum & -- \\
\hline

35--40 & Sicherung & Schreibt häufige Fehler an Tafel, korrigiert gemeinsam & Notieren korrekte Formen & Plenum & Tafel \\
\hline

40--45 & Über\-leitung & Erklärt Aufgabe für 2. Stunde (Einkaufsdialog) & Hören zu, bilden Paare & Plenum & -- \\
\hline
\end{longtable}

\vspace{1em}

\textbf{2. Stunde (45 Min.) -- Kommunikative Anwendung}

\begin{longtable}{|p{1cm}|p{1.5cm}|p{4cm}|p{3cm}|p{2cm}|p{2cm}|}
\hline
\textbf{Zeit} & \textbf{Phase} & \textbf{Lehrerhandeln} & \textbf{Schülerhandeln} & \textbf{SF} & \textbf{Medien} \\
\hline
\endhead

0--5 & Warm-up & Schnelles Quiz: Zeigt Bild, fragt \textit{¿Cuánto cuesta?} & Schätzen Preise auf Spanisch & Plenum & Bildkarten \\
\hline

5--15 & Anwen\-dung I & Modelliert Einkaufsdialog mit Beispiel, verteilt Rollenkarten & Üben Dialog in Partnerarbeit & PA & Rollenkarten \\
\hline

15--25 & Anwen\-dung II & Bildet Gruppen für Mini-Modenschau, gibt Anleitung & Bereiten kurze Präsentation vor (3-4 Sätze pro Person) & GA & -- \\
\hline

25--35 & Präsen\-tation & Moderiert Mini-Modenschau, gibt Feedback & Präsentieren: \textit{María lleva...} & Plenum & -- \\
\hline

35--42 & Festigung & Verteilt AB, unterstützt bei Schwierigkeiten & Bearbeiten AB-05 (Aufgaben 1-4) & EA & AB-05 \\
\hline

42--45 & Abschluss & Zusammenfassung, HA: LB-05 lesen, Vokabeln üben & Notieren HA, packen zusammen & Plenum & -- \\
\hline
\end{longtable}

\textbf{Legende:} EA = Einzelarbeit, PA = Partnerarbeit, GA = Gruppenarbeit, SF = Sozialform

% ============================================================================
% 7. ERWARTETE SCHWIERIGKEITEN
% ============================================================================
\section{Erwartete Schwierigkeiten}

\begin{tabular}{|p{5cm}|p{8cm}|}
\hline
\textbf{Schwierigkeit} & \textbf{Reaktion} \\
\hline
Farbadjektiv-Kongruenz (*la camiseta rojo) & Explizite Übung: Farbe + Kleidung zuordnen, Merksatz: ``Die Farbe passt sich an!'' \\
\hline
Unveränderliche Farben (rosa, naranja) & Liste an Tafel, Merkspruch: ``Rosa bleibt rosa, egal was kommt.'' \\
\hline
Plural von Farben (azul $\rightarrow$ azules) & Kontrastübung: el zapato azul $\rightarrow$ los zapatos azules \\
\hline
Spaltungsverb \textit{costar} & Visualisierung: cuesta (Sg.) vs. cuestan (Pl.) \\
\hline
Zeitproblem & Puffer: Modenschau kürzen oder als HA \\
\hline
Heterogenität & [B]: Bildkarten als Hilfe; [E]: Modenschau moderieren \\
\hline
\end{tabular}

% ============================================================================
% 8. HAUSAUFGABE
% ============================================================================
\section{Hausaufgabe}

\begin{itemize}
    \item \textbf{LB-05-repaso-ropa.pdf} lesen und Vokabeln wiederholen
    \item Vokabeln üben: Kleidung + Farben (S. 88-89 im Lehrwerk)
    \item Optional: Eigenes Outfit fotografieren und auf Spanisch beschreiben
\end{itemize}

% ============================================================================
% 9. ANLAGEN
% ============================================================================
\section{Anlagen}

\begin{enumerate}
    \item \textbf{AB-05-repaso-ropa.pdf} -- Arbeitsblatt (24× kopieren, max. 2 Seiten)
    \item \textbf{ML-05-repaso-ropa.pdf} -- Musterlösung (1× für Lehrkraft)
    \item \textbf{LB-05-repaso-ropa.pdf} -- Lernblatt (24× kopieren, max. 2 Seiten)
    \item \textbf{LH-05-repaso-ropa.pdf} -- Lehrerhinweise
    \item Bingo-Karten Kleidung (Vorlage im Anhang)
    \item Bildkarten Kleidungsstücke (laminiert)
    \item Rollenkarten Einkaufsdialog
\end{enumerate}

% ============================================================================
% 10. DIDAKTIK-SCORE
% ============================================================================
\newpage
\section{Didaktik-Score}

\begin{scorebox}
\textbf{Bewertung der didaktischen Qualität nach 10 Dimensionen}\\
\textbf{Ziel-Score: $\geq$ 9.0 (Premium-Qualität)}

\vspace{1em}
\begin{tabular}{|c|l|c|c|p{4.5cm}|}
\hline
\textbf{\#} & \textbf{Dimension} & \textbf{Gew.} & \textbf{Score} & \textbf{Notiz} \\
\hline
1 & Differenzierung & 15\% & 9/10 & [B]/[S]/[E] klar unterschieden, qualitative Abstufung \\
\hline
2 & Taxonomie (AFB) & 10\% & 9/10 & I: 30\%, II: 50\%, III: 20\% \\
\hline
3 & Lernziel-Alignment & 10\% & 10/10 & Alle Ziele aus Rahmenplan MV \\
\hline
4 & Aktivierung & 10\% & 10/10 & Spielerisch: Bingo, Kettenübung, Modenschau \\
\hline
5 & Scaffolding & 10\% & 9/10 & Bildkarten, Rollenkarten, Modellierung \\
\hline
6 & Feedback-Qualität & 10\% & 9/10 & Fehlerkorrektur an Tafel, Peer-Feedback \\
\hline
7 & Sprachliche Klarheit & 10\% & 10/10 & Altersgerechte Formulierungen \\
\hline
8 & Motivation \& Relevanz & 10\% & 10/10 & Alltagsbezug: Mode, eigene Kleidung \\
\hline
9 & Inklusion (UDL) & 10\% & 9/10 & Visuell + auditiv + kinästhetisch \\
\hline
10 & Praktikabilität & 5\% & 9/10 & 90 Min. realistisch, Materialien einfach \\
\hline
\multicolumn{3}{|r|}{\textbf{GESAMT:}} & \textbf{9.3/10} & \textcolor{successgreen}{\textbf{FREIGABE}} \\
\hline
\end{tabular}

\vspace{1em}
\textbf{Bewertungsschwellen:}
\begin{itemize}
    \item[$\boxtimes$] \textcolor{successgreen}{$\geq$ 9.0: \textbf{FREIGABE} (Premium-Qualität)}
    \item[$\Box$] \textcolor{warningorange}{8.0--8.9: Iteration empfohlen}
    \item[$\Box$] 6.0--7.9: Überarbeitung erforderlich
    \item[$\Box$] $<$ 6.0: Grundlegende Neukonzeption
\end{itemize}
\end{scorebox}

\end{document}
